\chapter{Chain Complexes}
\label{ch:v22}

Chain complexes organize sequences of modules with consecutive compositions zero. Homology measures the failure of exactness and is the central invariant of homological algebra.

\section*{Learning goals}
After this chapter, the reader should be able to:
\begin{itemize}
\item verify that a sequence of modules and differentials forms a chain complex by checking consecutive composites vanish;
\item compute cycles, boundaries, and homology groups in explicit finite complexes;
\item construct chain maps and verify compatibility with differentials;
\item identify exact complexes by vanishing homology;
\item compute mapping behavior induced on homology by chain maps;
\item distinguish chain complexes, exact sequences, and graded modules by their additional differential structure;
\end{itemize}
\section*{Conceptual roadmap}
\volroadmap
  {presentation or universal property}
  {exactness/localization}
  {structural consequence}

\section{Chain complexes}
A chain complex has modules $C_n$ and differentials $d_n:C_n\to C_{n-1}$ with $d_{n-1}d_n=0$.

\section{Cycles and boundaries}
Cycles are kernels of differentials; boundaries are images of the next differentials.

% BEGIN VOL05-EXPANSION V22-example-01
\begin{example}[A two-term complex]\label{ex:v22-expand-01}
The sequence $0\to\mathbb Z\xrightarrow{\cdot2}\mathbb Z\to0$ is a chain complex. Its degree-zero homology is $\mathbb Z/2\mathbb Z$, while the preceding homology vanishes because multiplication by $2$ is injective.
\end{example}
% END VOL05-EXPANSION V22-example-01
\section{Homology}
$H_n(C)=\ker d_n/\operatorname{im}d_{n+1}$.

\section{Exact complexes}
A complex is exact exactly when all homology modules vanish.

\section{Chain maps}
A chain map commutes with differentials and therefore sends cycles to cycles and boundaries to boundaries.

% BEGIN VOL05-EXPANSION V22-example-02
\begin{example}[Boundary squared is zero]\label{ex:v22-expand-02}
In any chain complex, $d_{n-1}d_n=0$, so every boundary is automatically a cycle. Homology measures the quotient of cycles by those cycles already explained as boundaries.
\end{example}
% END VOL05-EXPANSION V22-example-02
\section{Induced maps}
Every chain map induces maps on homology.

\section{Short exact sequences of complexes}
Degreewise short exact sequences produce long exact sequences in homology.

% BEGIN VOL05-EXPANSION V22-example-03
\begin{example}[Exact complex]\label{ex:v22-expand-03}
For a short exact sequence $0\to A\to B\to C\to0$ regarded as a complex, every homology group vanishes. Exactness is precisely acyclicity for this finite complex.
\end{example}
% END VOL05-EXPANSION V22-example-03
\section{Homotopies}
Chain-homotopic maps induce the same map on homology.

\section{Examples}
Two-term complexes, resolutions, Koszul-style complexes, and singular-like algebraic complexes illustrate the formalism.

% VOL05-EXACT-COMPUTATION-CHAINS-V22
\begin{remark}[Exact chain identities versus numerical residuals]
For matrices representing differentials over an exact coefficient ring, the
chain condition
\[
d_{n-1}d_n=0
\]
is an exact algebraic identity.  Likewise,
\[
H_n(C)=\ker d_n/\operatorname{im}d_{n+1}
\]
requires exact kernels, images, and containment of boundaries in cycles.

If a complex is entered using floating-point approximations to coefficients,
a small norm
\[
\|D_{n-1}D_n\|
\]
is only a numerical residual; it does not make the rounded matrices an exact
chain complex over the intended ring.  Conversely, an exact symbolic product
equal to zero is stronger than a tiny residual and should be verified exactly
whenever the input data are exact.

Over a field, exact row reduction may compute kernels and images.  Over
\(\mathbb Z\) one instead needs arithmetic respecting the integral module
structure, such as unimodular transformations or Smith-type methods; over
polynomial rings, syzygy or Gr\"obner-basis methods may be appropriate.
The algorithm must match the coefficient category in which the homology is
being claimed.
\end{remark}

\section{Core structural results}

\begin{theorem}[Chain maps induce homology maps]\label{thm:v22-01}
A chain map $f:C\to D$ induces $H_n(f):H_n(C)\to H_n(D)$.
\begin{proof}
Commutation with differentials sends cycles to cycles. It also sends boundaries to boundaries because $f_n d_{n+1}=d_{n+1}f_{n+1}$, so the map descends to the quotient.
\end{proof}
\end{theorem}

\begin{theorem}[Exactness equals zero homology]\label{thm:v22-02}
A chain complex is exact at every degree iff $H_n(C)=0$ for all $n$.
\begin{proof}
By definition, zero homology at degree $n$ means the cycle module equals the boundary module, which is exactly exactness at that term.
\end{proof}
\end{theorem}

\begin{theorem}[Chain homotopic maps agree on homology]\label{thm:v22-03}
If $f-g=dh+hd$, then $f$ and $g$ induce the same homology maps.
\begin{proof}
For a cycle $z$, $(f-g)(z)=d h(z)$ because $d(z)=0$, so the difference of images is a boundary.
\end{proof}
\end{theorem}

\section{Worked examples}

\begin{example}[Chain complexes diagnostic]\label{ex:v22-worked-01}
Analyze Chain complexes in the setting of this chapter: state the precise algebraic hypotheses, verify \(d^2=0\), then compute kernels, images, cycles, boundaries, and homology, and derive the resulting conclusion. A chain complex has modules $C_n$ and differentials $d_n:C_n\to C_{n-1}$ with $d_{n-1}d_n=0$. The chapter-specific mechanism is to verify \(d^2=0\), then compute kernels, images, cycles, boundaries, and homology; the hypotheses determine exactly which algebraic conclusion is available.
\end{example}

\begin{example}[Cycles and boundaries diagnostic]\label{ex:v22-worked-02}
Analyze Cycles and boundaries in the setting of this chapter: state the precise algebraic hypotheses, verify \(d^2=0\), then compute kernels, images, cycles, boundaries, and homology, and derive the resulting conclusion. Cycles are kernels of differentials; boundaries are images of the next differentials. The chapter-specific mechanism is to verify \(d^2=0\), then compute kernels, images, cycles, boundaries, and homology; the hypotheses determine exactly which algebraic conclusion is available.
\end{example}

\begin{example}[Homology diagnostic]\label{ex:v22-worked-03}
Analyze Homology in the setting of this chapter: state the precise algebraic hypotheses, verify \(d^2=0\), then compute kernels, images, cycles, boundaries, and homology, and derive the resulting conclusion. $H_n(C)=\ker d_n/\operatorname{im}d_{n+1}$. The chapter-specific mechanism is to verify \(d^2=0\), then compute kernels, images, cycles, boundaries, and homology; the hypotheses determine exactly which algebraic conclusion is available.
\end{example}

\begin{example}[Exact complexes diagnostic]\label{ex:v22-worked-04}
Analyze Exact complexes in the setting of this chapter: state the precise algebraic hypotheses, verify \(d^2=0\), then compute kernels, images, cycles, boundaries, and homology, and derive the resulting conclusion. A complex is exact exactly when all homology modules vanish. The chapter-specific mechanism is to verify \(d^2=0\), then compute kernels, images, cycles, boundaries, and homology; the hypotheses determine exactly which algebraic conclusion is available.
\end{example}

\section{Solved dossiers}
Each dossier is a canonical solved problem. The provenance ledger records which retained corpus rules guided its topic and which dossiers were added freshly to complete the chapter.

\begin{problem}[Chain complexes diagnostic]\label{prob:v22-dossier-01}
Analyze Chain complexes in the setting of this chapter: state the precise algebraic hypotheses, verify \(d^2=0\), then compute kernels, images, cycles, boundaries, and homology, and derive the resulting conclusion.
\end{problem}
\begin{solution}
A chain complex has modules $C_n$ and differentials $d_n:C_n\to C_{n-1}$ with $d_{n-1}d_n=0$. The chapter-specific mechanism is to verify \(d^2=0\), then compute kernels, images, cycles, boundaries, and homology; the hypotheses determine exactly which algebraic conclusion is available.
\end{solution}

\begin{problem}[Cycles and boundaries diagnostic]\label{prob:v22-dossier-02}
Analyze Cycles and boundaries in the setting of this chapter: state the precise algebraic hypotheses, verify \(d^2=0\), then compute kernels, images, cycles, boundaries, and homology, and derive the resulting conclusion.
\end{problem}
\begin{solution}
Cycles are kernels of differentials; boundaries are images of the next differentials. The chapter-specific mechanism is to verify \(d^2=0\), then compute kernels, images, cycles, boundaries, and homology; the hypotheses determine exactly which algebraic conclusion is available.
\end{solution}

\begin{problem}[Homology diagnostic]\label{prob:v22-dossier-03}
Analyze Homology in the setting of this chapter: state the precise algebraic hypotheses, verify \(d^2=0\), then compute kernels, images, cycles, boundaries, and homology, and derive the resulting conclusion.
\end{problem}
\begin{solution}
$H_n(C)=\ker d_n/\operatorname{im}d_{n+1}$. The chapter-specific mechanism is to verify \(d^2=0\), then compute kernels, images, cycles, boundaries, and homology; the hypotheses determine exactly which algebraic conclusion is available.
\end{solution}

\begin{problem}[Exact complexes diagnostic]\label{prob:v22-dossier-04}
Analyze Exact complexes in the setting of this chapter: state the precise algebraic hypotheses, verify \(d^2=0\), then compute kernels, images, cycles, boundaries, and homology, and derive the resulting conclusion.
\end{problem}
\begin{solution}
A complex is exact exactly when all homology modules vanish. The chapter-specific mechanism is to verify \(d^2=0\), then compute kernels, images, cycles, boundaries, and homology; the hypotheses determine exactly which algebraic conclusion is available.
\end{solution}

\begin{problem}[Chain maps diagnostic]\label{prob:v22-dossier-05}
Analyze Chain maps in the setting of this chapter: state the precise algebraic hypotheses, verify \(d^2=0\), then compute kernels, images, cycles, boundaries, and homology, and derive the resulting conclusion.
\end{problem}
\begin{solution}
A chain map commutes with differentials and therefore sends cycles to cycles and boundaries to boundaries. The chapter-specific mechanism is to verify \(d^2=0\), then compute kernels, images, cycles, boundaries, and homology; the hypotheses determine exactly which algebraic conclusion is available.
\end{solution}

\begin{problem}[Induced maps diagnostic]\label{prob:v22-dossier-06}
Analyze Induced maps in the setting of this chapter: state the precise algebraic hypotheses, verify \(d^2=0\), then compute kernels, images, cycles, boundaries, and homology, and derive the resulting conclusion.
\end{problem}
\begin{solution}
Every chain map induces maps on homology. The chapter-specific mechanism is to verify \(d^2=0\), then compute kernels, images, cycles, boundaries, and homology; the hypotheses determine exactly which algebraic conclusion is available.
\end{solution}

\begin{problem}[Short exact sequences of complexes diagnostic]\label{prob:v22-dossier-07}
Analyze Short exact sequences of complexes in the setting of this chapter: state the precise algebraic hypotheses, verify \(d^2=0\), then compute kernels, images, cycles, boundaries, and homology, and derive the resulting conclusion.
\end{problem}
\begin{solution}
Degreewise short exact sequences produce long exact sequences in homology. The chapter-specific mechanism is to verify \(d^2=0\), then compute kernels, images, cycles, boundaries, and homology; the hypotheses determine exactly which algebraic conclusion is available.
\end{solution}

\begin{problem}[Homotopies diagnostic]\label{prob:v22-dossier-08}
Analyze Homotopies in the setting of this chapter: state the precise algebraic hypotheses, verify \(d^2=0\), then compute kernels, images, cycles, boundaries, and homology, and derive the resulting conclusion.
\end{problem}
\begin{solution}
Chain-homotopic maps induce the same map on homology. The chapter-specific mechanism is to verify \(d^2=0\), then compute kernels, images, cycles, boundaries, and homology; the hypotheses determine exactly which algebraic conclusion is available.
\end{solution}

\begin{problem}[Chain maps induce homology maps proof dossier]\label{prob:v22-dossier-09}
Prove the following structural statement and identify the hypothesis that makes the conclusion work: A chain map $f:C\to D$ induces $H_n(f):H_n(C)\to H_n(D)$.
\end{problem}
\begin{solution}
Commutation with differentials sends cycles to cycles. It also sends boundaries to boundaries because $f_n d_{n+1}=d_{n+1}f_{n+1}$, so the map descends to the quotient. The reusable algebraic mechanism here is to verify \(d^2=0\), then compute kernels, images, cycles, boundaries, and homology.
\end{solution}

\begin{problem}[Exactness equals zero homology proof dossier]\label{prob:v22-dossier-10}
Prove the following structural statement and identify the hypothesis that makes the conclusion work: A chain complex is exact at every degree iff $H_n(C)=0$ for all $n$.
\end{problem}
\begin{solution}
By definition, zero homology at degree $n$ means the cycle module equals the boundary module, which is exactly exactness at that term. The reusable algebraic mechanism here is to verify \(d^2=0\), then compute kernels, images, cycles, boundaries, and homology.
\end{solution}

\begin{problem}[Chain homotopic maps agree on homology proof dossier]\label{prob:v22-dossier-11}
Prove the following structural statement and identify the hypothesis that makes the conclusion work: If $f-g=dh+hd$, then $f$ and $g$ induce the same homology maps.
\end{problem}
\begin{solution}
For a cycle $z$, $(f-g)(z)=d h(z)$ because $d(z)=0$, so the difference of images is a boundary. The reusable algebraic mechanism here is to verify \(d^2=0\), then compute kernels, images, cycles, boundaries, and homology.
\end{solution}

\begin{problem}[Chapter synthesis]\label{prob:v22-dossier-12}
Connect the definitions and principal structural theorems of this chapter into one reusable commutative-algebra strategy.
\end{problem}
\begin{solution}
Chain complexes organize sequences of modules with consecutive compositions zero. Homology measures the failure of exactness and is the central invariant of homological algebra. A chapter-specific strategy is to verify \(d^2=0\), then compute kernels, images, cycles, boundaries, and homology; consequently, chain complexes encode compatibility conditions and homology measures exactness failure.
\end{solution}

\section{Exercises with complete solutions}

\begin{exercise}\label{exr:v22-01}
Use Chain complexes to carry out the chapter's main algebraic method: verify \(d^2=0\), then compute kernels, images, cycles, boundaries, and homology. State one concrete consequence that follows under the stated hypotheses.
\end{exercise}
\begin{hint}
Check \(d_n d_{n+1}=0\) at every degree before calling the sequence a chain complex.
\end{hint}
\begin{solution}
A chain complex has modules $C_n$ and differentials $d_n:C_n\to C_{n-1}$ with $d_{n-1}d_n=0$. In this chapter the concrete consequence is that chain complexes encode compatibility conditions and homology measures exactness failure.
\end{solution}

\begin{exercise}\label{exr:v22-02}
Use Cycles and boundaries to carry out the chapter's main algebraic method: verify \(d^2=0\), then compute kernels, images, cycles, boundaries, and homology. State one concrete consequence that follows under the stated hypotheses.
\end{exercise}
\begin{hint}
Cycles are \(\ker d_n\) and boundaries are \(\operatorname{im}d_{n+1}\); homology is \(Z_n/B_n\).
\end{hint}
\begin{solution}
Cycles are kernels of differentials; boundaries are images of the next differentials. In this chapter the concrete consequence is that chain complexes encode compatibility conditions and homology measures exactness failure.
\end{solution}

\begin{exercise}\label{exr:v22-03}
Use Homology to carry out the chapter's main algebraic method: verify \(d^2=0\), then compute kernels, images, cycles, boundaries, and homology. State one concrete consequence that follows under the stated hypotheses.
\end{exercise}
\begin{hint}
To prove exactness at degree n, show every cycle is already a boundary.
\end{hint}
\begin{solution}
$H_n(C)=\ker d_n/\operatorname{im}d_{n+1}$. In this chapter the concrete consequence is that chain complexes encode compatibility conditions and homology measures exactness failure.
\end{solution}

\begin{exercise}\label{exr:v22-04}
Use Exact complexes to carry out the chapter's main algebraic method: verify \(d^2=0\), then compute kernels, images, cycles, boundaries, and homology. State one concrete consequence that follows under the stated hypotheses.
\end{exercise}
\begin{hint}
A chain map f satisfies d f = f d degreewise; write the commuting square at the degree in question.
\end{hint}
\begin{solution}
A complex is exact exactly when all homology modules vanish. In this chapter the concrete consequence is that chain complexes encode compatibility conditions and homology measures exactness failure.
\end{solution}

\begin{exercise}\label{exr:v22-05}
Use Chain maps to carry out the chapter's main algebraic method: verify \(d^2=0\), then compute kernels, images, cycles, boundaries, and homology. State one concrete consequence that follows under the stated hypotheses.
\end{exercise}
\begin{hint}
An induced homology map is well-defined because chain maps send cycles to cycles and boundaries to boundaries.
\end{hint}
\begin{solution}
A chain map commutes with differentials and therefore sends cycles to cycles and boundaries to boundaries. In this chapter the concrete consequence is that chain complexes encode compatibility conditions and homology measures exactness failure.
\end{solution}

\begin{exercise}\label{exr:v22-06}
Use Induced maps to carry out the chapter's main algebraic method: verify \(d^2=0\), then compute kernels, images, cycles, boundaries, and homology. State one concrete consequence that follows under the stated hypotheses.
\end{exercise}
\begin{hint}
For a short complex, compute kernels and images explicitly rather than using homological terminology as a substitute for calculation.
\end{hint}
\begin{solution}
Every chain map induces maps on homology. In this chapter the concrete consequence is that chain complexes encode compatibility conditions and homology measures exactness failure.
\end{solution}

\begin{exercise}\label{exr:v22-07}
Use Short exact sequences of complexes to carry out the chapter's main algebraic method: verify \(d^2=0\), then compute kernels, images, cycles, boundaries, and homology. State one concrete consequence that follows under the stated hypotheses.
\end{exercise}
\begin{hint}
Degree shifts and sign conventions matter when comparing chain and cochain complexes; state the convention first.
\end{hint}
\begin{solution}
Degreewise short exact sequences produce long exact sequences in homology. In this chapter the concrete consequence is that chain complexes encode compatibility conditions and homology measures exactness failure.
\end{solution}

\begin{exercise}\label{exr:v22-08}
Use Homotopies to carry out the chapter's main algebraic method: verify \(d^2=0\), then compute kernels, images, cycles, boundaries, and homology. State one concrete consequence that follows under the stated hypotheses.
\end{exercise}
\begin{hint}
Zero homology at every degree means exactness, but an exact complex need not be split or contractible.
\end{hint}
\begin{solution}
Chain-homotopic maps induce the same map on homology. In this chapter the concrete consequence is that chain complexes encode compatibility conditions and homology measures exactness failure.
\end{solution}

\section*{Chapter summary}
The chapter now contributes a canonical solved layer to the progression from rings and modules through localization, Noetherian methods, integral dependence, and derived functors.
% BEGIN VOL05-EXPANSION V22-exercises-01
\section{Graded supplementary exercises}
These exercises extend the protected chapter with a balanced set of computations, proofs, hypothesis tests, applications, and challenges.

\subsection*{Standard computations}

\begin{exercise}[Homology]\label{exr:v22-expand-01}
Compute the cokernel homology of $\mathbb Z\xrightarrow{\cdot3}\mathbb Z$.
\end{exercise}
\begin{hint}
Use cycles, boundaries, and homology.
\end{hint}
\begin{solution}
It is $\mathbb Z/3$.
\end{solution}

\begin{exercise}[Kernel homology]\label{exr:v22-expand-02}
Compute the kernel of multiplication by $3$ on $\mathbb Z$.
\end{exercise}
\begin{hint}
Use cycles, boundaries, and homology.
\end{hint}
\begin{solution}
It is zero.
\end{solution}

\begin{exercise}[Zero differential]\label{exr:v22-expand-03}
If all differentials are zero, what is homology?
\end{exercise}
\begin{hint}
Use cycles, boundaries, and homology.
\end{hint}
\begin{solution}
It equals the chain modules themselves.
\end{solution}

\begin{exercise}[Exact complex]\label{exr:v22-expand-04}
What is the homology of an exact complex?
\end{exercise}
\begin{hint}
Use cycles, boundaries, and homology.
\end{hint}
\begin{solution}
It is zero in every degree.
\end{solution}

\begin{exercise}[Boundary inclusion]\label{exr:v22-expand-05}
Why is every boundary a cycle?
\end{exercise}
\begin{hint}
Use cycles, boundaries, and homology.
\end{hint}
\begin{solution}
Because consecutive differentials compose to zero.
\end{solution}

\subsection*{Proofs}

\begin{exercise}[Chain maps induce homology maps proof]\label{exr:v22-expand-06}
Prove: A chain map $f:C\to D$ induces $H_n(f):H_n(C)\to H_n(D)$.
\end{exercise}
\begin{hint}
Start from the theorem hypotheses and identify the decisive algebraic construction.
\end{hint}
\begin{solution}
Commutation with differentials sends cycles to cycles. It also sends boundaries to boundaries because $f_n d_{n+1}=d_{n+1}f_{n+1}$, so the map descends to the quotient.
\end{solution}

\begin{exercise}[Exactness equals zero homology proof]\label{exr:v22-expand-07}
Prove: A chain complex is exact at every degree iff $H_n(C)=0$ for all $n$.
\end{exercise}
\begin{hint}
Start from the theorem hypotheses and identify the decisive algebraic construction.
\end{hint}
\begin{solution}
By definition, zero homology at degree $n$ means the cycle module equals the boundary module, which is exactly exactness at that term.
\end{solution}

\begin{exercise}[Chain homotopic maps agree on homology proof]\label{exr:v22-expand-08}
Prove: If $f-g=dh+hd$, then $f$ and $g$ induce the same homology maps.
\end{exercise}
\begin{hint}
Start from the theorem hypotheses and identify the decisive algebraic construction.
\end{hint}
\begin{solution}
For a cycle $z$, $(f-g)(z)=d h(z)$ because $d(z)=0$, so the difference of images is a boundary.
\end{solution}

\begin{exercise}[Structural synthesis]\label{exr:v22-expand-09}
Prove a corollary that genuinely uses both Chain maps induce homology maps and Exactness equals zero homology.
\end{exercise}
\begin{hint}
Apply the first theorem to obtain its structural description, then verify the second theorem's hypotheses.
\end{hint}
\begin{solution}
A correct proof explicitly invokes both results, verifies the common hypotheses, and derives a new consequence rather than merely restating either theorem.
\end{solution}

\subsection*{Counterexamples and hypothesis tests}

\begin{exercise}[Arbitrary maps]\label{exr:v22-expand-10}
Test the claim: Any sequence of module maps is a chain complex.
\end{exercise}
\begin{hint}
Use the smallest concrete ring, module, or resolution that can decide the claim.
\end{hint}
\begin{solution}
False: consecutive maps must compose to zero.
\end{solution}

\begin{exercise}[Cycles boundaries]\label{exr:v22-expand-11}
Test the claim: Every cycle is a boundary in every complex.
\end{exercise}
\begin{hint}
Use the smallest concrete ring, module, or resolution that can decide the claim.
\end{hint}
\begin{solution}
False; the quotient measures homology.
\end{solution}

\begin{exercise}[Zero homology]\label{exr:v22-expand-12}
Test the claim: Vanishing homology at one degree implies the whole complex is exact.
\end{exercise}
\begin{hint}
Use the smallest concrete ring, module, or resolution that can decide the claim.
\end{hint}
\begin{solution}
False; exactness must hold at every degree.
\end{solution}

\subsection*{Applications and investigations}

\begin{exercise}[Obstruction measure]\label{exr:v22-expand-13}
What does homology measure?
\end{exercise}
\begin{hint}
Translate the question into cycles, boundaries, and homology.
\end{hint}
\begin{solution}
It measures failure of exactness at each degree.
\end{solution}

\begin{exercise}[Topological bridge]\label{exr:v22-expand-14}
Why are chain complexes useful beyond algebra?
\end{exercise}
\begin{hint}
Translate the question into cycles, boundaries, and homology.
\end{hint}
\begin{solution}
They encode boundary operators whose homology captures invariants of spaces and resolutions.
\end{solution}

\subsection*{Challenge problems}

\begin{exercise}[Local-global synthesis]\label{exr:v22-expand-15}
Connect 'Chain complexes' with 'Examples' in one rigorous argument.
\end{exercise}
\begin{hint}
State the relevant ring map, module map, or complex before applying the structural theorem.
\end{hint}
\begin{solution}
The opening idea is: A chain complex has modules $C_n$ and differentials $d_n:C_n\to C_{n-1}$ with $d_{n-1}d_n=0$. The later viewpoint is: Two-term complexes, resolutions, Koszul-style complexes, and singular-like algebraic complexes illustrate the formalism. A complete solution identifies the structural theorem that links these descriptions and checks its hypotheses.
\end{solution}

\begin{exercise}[Hypothesis audit]\label{exr:v22-expand-16}
Choose one theorem from the chapter, remove one essential hypothesis, and exhibit or explain a concrete failure.
\end{exercise}
\begin{hint}
Use one of the chapter's explicit computations or counterexamples as the test case.
\end{hint}
\begin{solution}
The solution must name the removed hypothesis, show exactly which proof step breaks, and give a concrete algebraic object where the claimed conclusion fails.
\end{solution}

% END VOL05-EXPANSION V22-exercises-01
